% GENERATED by tools/coverage.py from the lesson plans -- do not edit by hand.
\subsection{Objectives per goal and week (generated)}
Each number is how many of that week's \enquote{The student can\ldots} objectives
serve the goal. Blank = not a focus that week (it may still be practised).

{\small\setlength{\tabcolsep}{3pt}
\begin{longtable}{@{}lcccccccccccccc@{}}
\toprule
Goal & 1 & 2 & 3 & 4 & 5 & 6 & 7 & 8 & 9 & 10 & 11 & 12 & B1 & B2\\ \midrule\endhead
L1.G1 & 2 &  &  &  &  & 1 & 1 &  &  &  & 1 & 1 & 2 & 1\\
L1.G2 &  & 3 & 1 &  &  & 1 &  & 3 &  &  & 1 &  &  & 1\\
L1.G3 &  &  & 2 & 1 & 1 & 1 &  &  & 1 &  & 1 &  &  & \\
L1.G4 &  &  &  & 2 &  & 1 &  &  & 2 & 1 & 1 &  &  & \\
L1.G5 &  &  &  &  & 2 & 1 &  &  &  & 2 & 1 &  &  & \\
L1.G6 & 1 & 1 & 1 & 1 & 1 & 1 & 1 & 1 &  & 1 & 1 &  & 1 & 1\\
L1.G7 & 1 &  & 1 &  &  & 1 &  & 1 &  &  & 1 & 2 & 1 & 1\\
L1.G8 & 1 &  & 1 & 1 & 1 & 1 & 1 &  & 1 &  &  & 2 &  & 1\\
L1.G9 &  &  &  &  &  & 1 & 2 & 1 & 1 &  & 1 & 2 &  & 1\\
\bottomrule\end{longtable}}

\subsection{All lesson objectives (generated)}
\begin{longtable}{@{}p{1.1cm}p{0.78\linewidth}p{1.6cm}@{}}
\toprule Week & Objective & Goals\\ \midrule\endhead
\multicolumn{3}{@{}l}{\textbf{1: Robots Need Instructions}}\\
1.1 & The student can write an algorithm as exact steps in order so that a \enquote{robot} draws the right picture. & L1.G1\\
1.2 & The student can point to the stage, a sprite, the blocks, the code area and the green flag in the Scratch app, and run a script. & L1.G1\\
1.3 & The student can find and fix a script whose steps are in the wrong order. & L1.G6\\
1.4 & The student can explain our class agreements and why we use a nickname and keep projects on our own Mac. & L1.G7, L1.G8\\
\midrule
\multicolumn{3}{@{}l}{\textbf{2: Loop-de-Loop}}\\
2.1 & The student can replace steps that repeat with a \enquote{repeat} block, and use \enquote{forever} to keep something going. & L1.G2\\
2.2 & The student can animate a sprite by switching costumes inside a loop. & L1.G2\\
2.3 & The student can add the Pen extension and draw a regular shape with \enquote{repeat n} and a turn of 360 ÷ n degrees. & L1.G2\\
2.4 & The student can find and fix a shape that does not close because of a wrong turn angle. & L1.G6\\
\midrule
\multicolumn{3}{@{}l}{\textbf{3: When This Happens\ldots}}\\
3.1 & The student can start different scripts with different events: green flag, a key press or a click on a sprite. & L1.G3\\
3.2 & The student can use \enquote{if} and \enquote{if-else} with a sensing question (\enquote{touching \ldots?}) to choose what a sprite does. & L1.G3\\
3.3 & The student can explain why an \enquote{if} must sit inside \enquote{forever} to keep checking, and fix a sprite that never reacts. & L1.G2, L1.G6\\
3.4 & The student can give a partner kind, specific feedback with \enquote{two stars and a wish}. & L1.G7, L1.G8\\
\midrule
\multicolumn{3}{@{}l}{\textbf{4: Keeping Score}}\\
4.1 & The student can make a variable, set it at the start of a game and change it while the game runs. & L1.G4\\
4.2 & The student can use \enquote{pick random} to put a sprite in a surprising place. & L1.G4\\
4.3 & The student can combine events, a forever loop and if-blocks to build a catch game, and fix a score that never resets or goes up too fast. & L1.G3, L1.G6\\
4.4 & The student can say whose art and sounds a project uses and add a credits note. & L1.G8\\
\midrule
\multicolumn{3}{@{}l}{\textbf{5: Messages and My Blocks}}\\
5.1 & The student can break a big job into smaller named parts (decomposition). & L1.G5\\
5.2 & The student can make a My Block with an input, like \enquote{polygon (sides) (length)}, and use it with different values. & L1.G5\\
5.3 & The student can make sprites work together with \enquote{broadcast} and \enquote{when I receive}, and fix a message that nobody receives. & L1.G3, L1.G6\\
5.4 & The student can choose colour-blind-safe colours, add shapes or words so colour is not the only clue, and keep text readable. & L1.G8\\
\midrule
\multicolumn{3}{@{}l}{\textbf{6: Review Studio and Checkpoint}}\\
6.1 & The student can match Scratch blocks to the Python lines that do the same job. & L1.G9\\
6.2 & The student can use the four detective steps to find and fix bugs in projects someone else made. & L1.G6, L1.G7\\
6.3 & The student can show what they know about sequence, loops, events, conditionals, variables and My Blocks in the mid-level checkpoint. & L1.G1, L1.G2, L1.G3, L1.G4, L1.G5\\
6.4 & The student can explain what makes a project fun and fair for everyone. & L1.G8\\
\midrule
\multicolumn{3}{@{}l}{\textbf{7: Hello, Python!}}\\
7.1 & The student can open a file in Thonny, press Run, find the Editor, the Shell and the turtle window, and save with a sensible file name & L1.G8, L1.G9\\
7.2 & The student can write \texttt{print} lines and comments, and explain that Python skips comments & L1.G9\\
7.3 & The student can turn Scratch \enquote{move} and \enquote{turn} blocks into turtle code that draws a square & L1.G1\\
7.4 & The student can fix a typo by reading the line number and the last line of a red error message & L1.G6\\
\midrule
\multicolumn{3}{@{}l}{\textbf{8: Loops in Python}}\\
8.1 & The student can write a \texttt{for i in range(n):} loop and say how many times it repeats & L1.G2, L1.G9\\
8.2 & The student can use indentation to show which lines are inside a loop & L1.G2\\
8.3 & The student can draw any regular polygon using the rule \enquote{turn = 360 ÷ sides}, and a pattern with a loop inside a loop & L1.G2\\
8.4 & The student can find and fix an indentation bug and an off-by-one bug, and ask for help with \enquote{I expected … but …} & L1.G6, L1.G7\\
\midrule
\multicolumn{3}{@{}l}{\textbf{9: Variables, Input and Choices}}\\
9.1 & The student can store a value in a variable and use it later in the program & L1.G4\\
9.2 & The student can ask a question with \texttt{input} and turn the answer into a number with \texttt{int} & L1.G4, L1.G9\\
9.3 & The student can use \texttt{if}, \texttt{elif} and \texttt{else} with comparisons (\texttt{==}, \texttt{<}, \texttt{>}) to choose what happens & L1.G3\\
9.4 & The student can explain which questions a program should never ask & L1.G8\\
\midrule
\multicolumn{3}{@{}l}{\textbf{10: Random Art and My Own Commands}}\\
10.1 & The student can define a function with \texttt{def} and parameters, and call it several times with different values & L1.G5\\
10.2 & The student can use \texttt{random.randint} and \texttt{random.choice} to make a picture that changes every run & L1.G4\\
10.3 & The student can find the bug \enquote{function defined but never called} & L1.G6\\
10.4 & The student can plan a project by sketching it and splitting it into parts & L1.G5\\
\midrule
\multicolumn{3}{@{}l}{\textbf{11: Project Studio}}\\
11.1 & The student can build a program in small steps: change one thing, run, check, repeat & L1.G1, L1.G9\\
11.2 & The student can make their project's main function work and call it in at least two places (Milestone M2) & L1.G5\\
11.3 & The student can add a question, an \texttt{if}, a loop and a random surprise to their project (Milestone M3) & L1.G2, L1.G3, L1.G4\\
11.4 & The student can test a partner's program, describe a bug kindly and act on feedback & L1.G6, L1.G7\\
\midrule
\multicolumn{3}{@{}l}{\textbf{12: Showcase and Celebration}}\\
12.1 & The student can run their project for a visitor and explain one part of the code in their own words (Milestone M4) & L1.G7, L1.G9\\
12.2 & The student can give and receive \enquote{two stars and a wish} feedback during a gallery walk & L1.G7, L1.G8\\
12.3 & The student can read, predict and fix short Python programs in the exit assessment & L1.G1, L1.G9\\
12.4 & The student can share work safely: on screen only, nicknames only, no personal information & L1.G8\\
\midrule
\multicolumn{3}{@{}l}{\textbf{B1: Unplugged Games Day}}\\
B1.1 & The student can turn a code of 0s and 1s (or a number code) into a pixel picture, and turn their own picture into a code. & L1.G1\\
B1.2 & The student can write and follow an exact algorithm as part of a team relay. & L1.G1, L1.G7\\
B1.3 & The student can use the four detective steps to find and fix the bug in an algorithm on a card. & L1.G6\\
\midrule
\multicolumn{3}{@{}l}{\textbf{B2: Free-Build and Remix Studio}}\\
B2.1 & The student can choose an earlier program and change it in at least two deliberate ways & L1.G1, L1.G2\\
B2.2 & The student can save a remix as a new file so the original stays safe & L1.G9\\
B2.3 & The student can give credit for ideas and code they reused & L1.G8\\
B2.4 & The student can find a bug a partner hid in a program & L1.G6, L1.G7\\
\midrule
\bottomrule\end{longtable}
